\chapter{Goitre}

\section*{Definition and Overview}
A goitre is an abnormal enlargement of the thyroid gland, the butterfly-shaped gland in the front of the neck that produces hormones controlling metabolism. Goitres are common and usually harmless, but they can be associated with thyroid dysfunction, causing overactivity (hyperthyroidism) or underactivity (hypothyroidism), or with nodules, some of which can be cancerous. The enlargement can be diffuse (affecting the whole gland) or nodular, and it can range from a barely noticeable bulge to a large mass that causes pressure symptoms such as difficulty swallowing or breathing. Treatment depends on the cause, the size, and whether the gland is functioning normally.

\section*{Epidemiology}
Goitre affects around 5--10\% of the population, and it is several times more common in women than in men. The most common cause worldwide is iodine deficiency, but in Australia, where salt and bread are fortified with iodine, autoimmune thyroid disease is the leading cause. Around 5--10\% of the Australian population has thyroid nodules, and about 10--15\% of people with a goitre have thyroid dysfunction. Goitres can develop at any age but become more common with increasing age. Most goitres are benign, with fewer than 5\% of solitary nodules being cancerous.

\section*{Aetiology and Pathophysiology}
A goitre develops when the thyroid gland enlarges in response to factors that stimulate thyroid growth or cause inflammation.

\begin{itemize}
    \item \textbf{Iodine deficiency:} the thyroid cannot make enough hormone without iodine, so the gland enlarges to compensate; now uncommon in Australia
    \item \textbf{Hashimoto thyroiditis:} autoimmune inflammation that gradually destroys the gland, often causing hypothyroidism and a goitre
    \item \textbf{Graves disease:} autoimmune stimulation of the thyroid causing overactivity and a diffuse goitre
    \item \textbf{Multinodular goitre:} age-related growth of multiple nodules, sometimes producing excess hormone
    \item \textbf{Solitary nodules:} benign cysts or adenomas, or rarely thyroid cancer
    \item \textbf{Other causes:} certain medicines, pregnancy, and inherited conditions
    \item \textbf{Thyroid-stimulating hormone (TSH):} when thyroid hormone levels fall, the pituitary produces more TSH, which stimulates the gland to grow
\end{itemize}

\section*{Clinical Features}
Many goitres cause no symptoms, but larger ones can cause pressure effects.

\begin{itemize}
    \item A visible swelling or lump at the front of the neck, below the Adam's apple
    \item A feeling of tightness or fullness in the throat
    \item Difficulty swallowing or a sensation of food sticking
    \item Hoarseness or voice changes
    \item Difficulty breathing, especially when lying flat, in large goitres
    \item Symptoms of hyperthyroidism: weight loss, palpitations, sweating, tremor, and anxiety
    \item Symptoms of hypothyroidism: fatigue, weight gain, cold intolerance, and constipation
    \item Cough or choking sensation in some cases
\end{itemize}

\section*{Differential Diagnosis}
Other causes of neck lumps and swelling should be considered.

\begin{itemize}
    \item Enlarged lymph nodes in the neck from infection or malignancy
    \item Thyroglossal cyst and branchial cyst
    \item Lipoma and other benign soft tissue lumps
    \item Salivary gland enlargement
    \item Submandibular gland stones or infection
    \item Carotid body tumour and other vascular lesions
    \item Thyroid cancer presenting as a nodule or goitre
\end{itemize}

\section*{When to Seek Medical Care}
See a doctor for any new lump or swelling in the neck, especially if it is growing, causing pressure symptoms, or associated with symptoms of thyroid dysfunction. A rapidly enlarging lump, hoarseness, or difficulty swallowing or breathing requires prompt assessment.

\textbf{Call 000 (or local emergency services) for:}
\begin{itemize}
    \item Sudden difficulty breathing or severe airway obstruction from a large goitre
    \item A rapidly enlarging neck lump with pain, fever, or difficulty swallowing
    \item Signs of thyroid storm: high fever, agitation, rapid heart rate, and confusion in a person with hyperthyroidism
\end{itemize}

\section*{Diagnosis and Investigations}
The cause and function of a goitre are determined with blood tests, imaging, and sometimes biopsy.

\begin{itemize}
    \item \textbf{Thyroid function tests:} TSH, free T4, and free T3 determine whether the gland is underactive, overactive, or normal
    \item \textbf{Thyroid antibodies:} detect autoimmune causes such as Hashimoto thyroiditis and Graves disease
    \item \textbf{Ultrasound:} assesses the size of the gland and characterises nodules
    \item \textbf{Thyroid scan:} a nuclear medicine test that shows whether nodules are hot (functioning) or cold (non-functioning)
    \item \textbf{Fine needle aspiration biopsy:} a sample from a suspicious nodule is examined for cancer cells
    \item \textbf{Iodine levels:} measured in areas where deficiency is suspected
\end{itemize}

\section*{Management}
Treatment depends on the cause, the size of the goitre, and the presence of symptoms or dysfunction.

\begin{itemize}
    \item \textbf{Observation:} small, asymptomatic goitres with normal function may simply be monitored with regular reviews and ultrasound
    \item \textbf{Iodine supplementation:} treats iodine deficiency where present
    \item \textbf{Thyroid hormone replacement:} levothyroxine corrects hypothyroidism and may slow further growth in some cases
    \item \textbf{Antithyroid medicines:} control hyperthyroidism in Graves disease
    \item \textbf{Radioactive iodine:} shrinks overactive thyroid tissue in hyperthyroidism
    \item \textbf{Surgery:} partial or total thyroidectomy for large goitres causing pressure, suspicious nodules, or cancer
    \item \textbf{Follow-up:} regular monitoring of function, size, and nodules
\end{itemize}

\section*{Complications}
\begin{itemize}
    \item Progression of the goitre causing airway or swallowing obstruction
    \item Thyroid dysfunction: hypothyroidism or hyperthyroidism with their systemic effects
    \item Bleeding into a nodule causing sudden swelling and pain
    \item Cancer in a nodule, which is uncommon but important to detect
    \item Complications of surgery, including damage to the voice box nerves, low calcium from parathyroid injury, and hypothyroidism
    \item Recurrence after partial removal
\end{itemize}

\section*{Prognosis and Outcomes}
The prognosis of goitre is generally excellent. Most goitres grow slowly and cause no problems, and many remain the same size for years. When dysfunction is present, it is easily corrected with medication. Surgery cures most symptomatic goitres, and thyroid cancer, when found, is usually treatable. The key to good outcomes is regular monitoring, because goitres can change over time, and new symptoms such as hoarseness or rapid growth should always be investigated.

\section*{Prevention and Self-Care}
\begin{itemize}
    \item Ensure adequate iodine intake through iodised salt, bread, and a varied diet
    \item Do not take high-dose iodine supplements unless advised, as excess can harm the thyroid
    \item Have thyroid function checked if you have a family history of thyroid disease or symptoms of dysfunction
    \item Examine your neck occasionally for swelling, and see a doctor if you notice a lump
    \item Attend recommended follow-up, including blood tests and ultrasound
    \item Take prescribed thyroid medicines exactly as directed
    \item Report hoarseness, difficulty swallowing, or rapid enlargement promptly
    \item Protect your neck during contact sports if you have a large goitre
\end{itemize}

\section*{Sources and Further Reading}
The following reputable sources provide further information for healthcare students and patients seeking reliable, current advice about goitre.

\begin{itemize}
    \item Australian Thyroid Foundation. \textit{Goitre and thyroid disorders information.} https://www.thyroidfoundation.org.au/
    \item Healthdirect Australia. \textit{Goitre.} https://www.healthdirect.gov.au/
    \item Royal Australasian College of Surgeons. \textit{Thyroid surgery patient information.} https://www.surgeons.org/
    \item NHS. \textit{Goitre: overview and treatment.} https://www.nhs.uk/
    \item Mayo Clinic. \textit{Goiter: symptoms and causes.} https://www.mayoclinic.org/
    \item MSD Manual. \textit{Simple goitre.} https://www.msdmanuals.com/
\end{itemize}
